\documentclass[11pt]{article}

\usepackage{amsmath,amssymb,amsthm,mathtools}
\usepackage[margin=1in]{geometry}
\usepackage{microtype}

\newtheorem{theorem}{Theorem}[section]
\newtheorem{lemma}[theorem]{Lemma}
\newtheorem{corollary}[theorem]{Corollary}
\newtheorem{remark}[theorem]{Remark}

\newcommand{\dd}{\,\mathrm d}
\newcommand{\one}{\mathbf 1}
\newcommand{\ip}[2]{\left\langle #1,#2\right\rangle}

\title{Triangle Books with a Two-Edge Tail:\\
A Sharp Operator-Square Inequality}
\author{}
\date{}

\begin{document}

\maketitle

\begin{abstract}
Let $R_m$ be a book of $m\geq1$ triangles on one common spine edge,
with one otherwise disjoint path of length two attached to a spine
endpoint.  We prove the chromatic-polynomial graphon bound for every
$R_m$ and every $p\geq1/2$.  The new input is a sharp inequality for the
weight $A=T_Wd$:
\[
 \int A\tau\geq(2p-1)\int Ad.
\]
Unlike a monotonicity assertion for the weight $A$, this follows from an
exact Hilbert-space square after writing $d=p\mathbf1+g$.  Moment convexity
under the positive measure $W(x,y)A(x)\,\mathrm d\mu(x)\mathrm d\mu(y)$
then treats all pages at once.  The case $m=2$ proves the formerly open
NetworkX Atlas graph $120$.
\end{abstract}


%%%%%%%%%%%%%%%%%%%%%%%%%%%%%%%%%%%%%%%%%%%%%%%%%%%%%%%%%%%%%%%%%%%%%%%%%%%%%%%
\section{The family and its exact target}
%%%%%%%%%%%%%%%%%%%%%%%%%%%%%%%%%%%%%%%%%%%%%%%%%%%%%%%%%%%%%%%%%%%%%%%%%%%%%%%

Let $(\Omega,\mu)$ be an arbitrary probability space and let
$W\colon\Omega^2\to[0,1]$ be a symmetric measurable graphon.  Write
\[
 T=T_W,\qquad
 p=\int_{\Omega^2}W\dd\mu^2,
 \qquad d=T\one,
 \qquad A=Td=T^2\one.
\]
We also use the rooted triangle density
\[
 \tau(x)=\int_{\Omega^2}
 W(x,y)W(x,z)W(y,z)\dd\mu(y)\dd\mu(z)
\]
and the common-neighbour kernel
\[
 S(x,y)=\int_\Omega W(x,z)W(y,z)\dd\mu(z).
\]
All these functions are measurable and take values in $[0,1]$.

For $m\geq1$, define $R_m$ on a spine $ab$, page vertices
$z_1,\ldots,z_m$, and tail vertices $u,v$.  Its edges are
\[
 ab,\quad az_i,bz_i\ (1\leq i\leq m),\quad au,uv.
\]
Thus $R_m$ is a book of $m$ triangles with one two-edge tail at the
spine endpoint $a$.  A proper colouring first colours the ordered spine,
then each page independently, and finally the two tail vertices.  Hence
\begin{equation}
 \chi_{R_m}(x)=x(x-1)^3(x-2)^m,
 \qquad \chi(R_m)=3.
 \label{eq:chromatic}
\end{equation}
Putting $q=1-p$ and substituting $x=1/q$ gives
\begin{equation}
 \Phi_{R_m}(p)
 =q^{m+4}\chi_{R_m}(1/q)
 =p^3(2p-1)^m.
 \label{eq:target}
\end{equation}
The formula is polynomial and therefore also gives the value at $p=1$.
The required interval is $p\in[1/2,1]$.


%%%%%%%%%%%%%%%%%%%%%%%%%%%%%%%%%%%%%%%%%%%%%%%%%%%%%%%%%%%%%%%%%%%%%%%%%%%%%%%
\section{A sharp inequality for the two-edge-tail weight}
%%%%%%%%%%%%%%%%%%%%%%%%%%%%%%%%%%%%%%%%%%%%%%%%%%%%%%%%%%%%%%%%%%%%%%%%%%%%%%%

The next lemma is the only new analytic ingredient.  It is important that
its proof does not insert $A$ into the previously established
degree-weighted rooted-Goodman lemma: $A$ is not a pointwise monotone
function of the complement degree, so that substitution would be invalid.

\begin{lemma}[$T_Wd$-weighted rooted triangle]
\label{lem:A-triangle}
If $p\geq1/6$, then
\begin{equation}
 \int_\Omega A(x)\tau(x)\dd\mu(x)
 \geq(2p-1)\int_\Omega A(x)d(x)\dd\mu(x).
 \label{eq:A-triangle}
\end{equation}
\end{lemma}

\begin{proof}
For $a,b\in[0,1]$ one has $ab\geq a+b-1$.  Apply this to
$W(x,y),W(x,z)$, multiply by the nonnegative $W(y,z)$, and integrate in
$y,z$.  Symmetry and Fubini give the pointwise estimate
\begin{equation}
 \tau(x)\geq2A(x)-p.
 \label{eq:pointwise}
\end{equation}
Since $A\geq0$, it follows that
\begin{equation}
 \int A\tau\geq2\lVert A\rVert_2^2-p\lVert d\rVert_2^2,
 \label{eq:first-reduction}
\end{equation}
where we used
\[
 \int A=\ip{Td}{\one}=\ip{d}{T\one}=\lVert d\rVert_2^2.
\]

It remains to prove
\begin{equation}
 E:=2\lVert A\rVert_2^2-p\lVert d\rVert_2^2
 -(2p-1)\ip{A}{d}\geq0.
 \label{eq:E}
\end{equation}
Put $g=d-p\one$ and $h=Tg$.  Then $\int g=0$ and
\[
 A=T(p\one+g)=pd+h.
\]
The coefficient of $\lVert d\rVert_2^2$ cancels exactly in
\eqref{eq:E}, leaving
\begin{equation}
 E=2\lVert h\rVert_2^2+(2p+1)\ip{h}{d}.
 \label{eq:E-one}
\end{equation}
Self-adjointness of $T$ and $T\one=d$ show that
\[
 \ip{h}{\one}=\ip{Tg}{\one}=\ip{g}{d}=\lVert g\rVert_2^2.
\]
Consequently
\[
 \ip{h}{d}
 =p\lVert g\rVert_2^2+\ip{h}{g}.
\]
Substitution in \eqref{eq:E-one} and completion of the square give the
exact identity
\begin{equation}
 E=
 2\left\lVert h+\frac{2p+1}{4}g\right\rVert_2^2
 +\frac{(2p+1)(6p-1)}8\lVert g\rVert_2^2.
 \label{eq:operator-square}
\end{equation}
Both coefficients are nonnegative for $p\geq1/6$.  Combining
\eqref{eq:first-reduction}--\eqref{eq:operator-square} proves
\eqref{eq:A-triangle}.
\end{proof}

\begin{remark}[Why the square is substantive]
No positivity of the graphon operator is assumed; a nonnegative symmetric
kernel can have negative eigenvalues.  Identity
\eqref{eq:operator-square} controls those spectral directions through a
quadratic polynomial with negative discriminant when $p>1/6$.  Thus the
proof is valid for arbitrary graphons, not only positive-semidefinite or
regular kernels.
\end{remark}


%%%%%%%%%%%%%%%%%%%%%%%%%%%%%%%%%%%%%%%%%%%%%%%%%%%%%%%%%%%%%%%%%%%%%%%%%%%%%%%
\section{The book-tail theorem}
%%%%%%%%%%%%%%%%%%%%%%%%%%%%%%%%%%%%%%%%%%%%%%%%%%%%%%%%%%%%%%%%%%%%%%%%%%%%%%%

Define
\[
 B=\int_\Omega A(x)d(x)\dd\mu(x),
 \qquad
 F=\int_\Omega A(x)\tau(x)\dd\mu(x).
\]
The quantity $B=t(P_4,W)$, while $F=t(R_1,W)$ is the density of a
triangle with one two-edge tail.

\begin{lemma}[Sharp lower bound for the first page]
\label{lem:first-page}
For $p\in[1/2,1]$,
\[
 F\geq p^3(2p-1).
\]
\end{lemma}

\begin{proof}
Let $M=\int A=\int d^2$.  From \eqref{eq:pointwise}, Cauchy--Schwarz,
and Jensen,
\[
 F\geq2\int A^2-p\int A\geq2M^2-pM.
\]
Also $M\geq p^2$.  The function $r\mapsto2r^2-pr$ is increasing for
$r\geq p^2$ when $p\geq1/2$, because there
$4r-p\geq4p^2-p\geq0$.  Hence
\[
 F\geq2p^4-p^3=p^3(2p-1).
\]
\end{proof}

\begin{theorem}[Triangle-book two-edge-tail theorem]
\label{thm:main}
Let $m\geq1$.  Every graphon $W$ of edge density
$p\in[1/2,1]$ satisfies
\[
 \boxed{
 t(R_m,W)\geq p^3(2p-1)^m=\Phi_{R_m}(p).
 }
\]
\end{theorem}

\begin{proof}
Because $p>0$, the forest inequality gives
$B=t(P_4,W)\geq p^3>0$.  Define a finite positive measure on
$\Omega^2$ by
\[
 \dd\nu(x,y)=W(x,y)A(x)\dd\mu(x)\dd\mu(y).
\]
Its total mass is
\[
 \nu(\Omega^2)=\int W(x,y)A(x)\dd\mu^2=B.
\]
Integrating all page vertices and the two tail vertices in the defining
homomorphism integral gives the exact identity
\begin{equation}
 t(R_m,W)=\int_{\Omega^2}S(x,y)^m\dd\nu(x,y).
 \label{eq:density}
\end{equation}
Moreover, symmetry and Fubini give
\begin{equation}
 \int_{\Omega^2}S(x,y)\dd\nu(x,y)
 =\int_\Omega A(x)\tau(x)\dd\mu(x)=F.
 \label{eq:first-moment}
\end{equation}
Apply Jensen's inequality to the convex function $s\mapsto s^m$ under
the probability measure $\nu/B$.  Equations
\eqref{eq:density}--\eqref{eq:first-moment} yield
\begin{equation}
 t(R_m,W)
 \geq B\left(\frac FB\right)^m
 =F\left(\frac FB\right)^{m-1}.
 \label{eq:moment}
\end{equation}
Lemma~\ref{lem:A-triangle} gives $F/B\geq2p-1$, and
Lemma~\ref{lem:first-page} gives $F\geq p^3(2p-1)$.  Since
$2p-1\geq0$, substitution in \eqref{eq:moment} proves
\[
 t(R_m,W)\geq p^3(2p-1)^m.
\]
This is the target in \eqref{eq:target}.
\end{proof}

At the threshold $p=1/2$ the target is zero; the proof above remains valid
and uses no division by $2p-1$.  At $p=1$, $W=1$ almost everywhere and
both sides equal one.  For every integer $k\geq2$, the balanced complete
$k$-partite graphon has $p=1-1/k$ and realizes equality: along a graphon
edge, $S(x,y)=1-2/k=2p-1$, while $d=p$ and $A=p^2$ almost everywhere.


%%%%%%%%%%%%%%%%%%%%%%%%%%%%%%%%%%%%%%%%%%%%%%%%%%%%%%%%%%%%%%%%%%%%%%%%%%%%%%%
\section{Atlas 120 and formalization boundary}
%%%%%%%%%%%%%%%%%%%%%%%%%%%%%%%%%%%%%%%%%%%%%%%%%%%%%%%%%%%%%%%%%%%%%%%%%%%%%%%

\begin{corollary}[Atlas 120]
NetworkX Atlas graph $120$, with graph6 code
\texttt{EB\char125?}, is positive.  It has order six, size seven,
chromatic number three, and
\[
 \chi_{F_{120}}(x)=x(x-1)^3(x-2)^2.
\]
Every graphon of edge density $p\in[1/2,1]$ satisfies
\[
 t(F_{120},W)\geq p^3(2p-1)^2=\Phi_{F_{120}}(p).
\]
\end{corollary}

\begin{proof}
Its NetworkX edge list is
\[
 04,\ 05,\ 13,\ 14,\ 23,\ 24,\ 34.
\]
Take spine $ab=43$, page vertices $1,2$, and tail $4-0-5$.
This is exactly $R_2$, so Theorem~\ref{thm:main} applies.
\end{proof}

Every use of Tonelli--Fubini above concerns bounded nonnegative measurable
integrands.  The operator calculation uses only bounded self-adjointness of
$T_W$ on $L^2$ and never assumes that $T_W$ is positive semidefinite.
The sole previously accepted density input is the forest bound for $P_4$,
used only to certify that the normalizing mass $B$ is positive; its full
sharp lower value is not substituted into the final estimate.

For formal verification, the reusable new lemma is
Lemma~\ref{lem:A-triangle}, including the pointwise triangle reduction and
the exact operator-square identity \eqref{eq:operator-square}.  The remaining
steps are the two Fubini identities, probability-measure normalization,
Jensen for integer moments, and the scalar first-page estimate.  The
graph-specific module awaiting inspection is
\texttt{lean/Taeyoung/Examples/Graph120.lean}.

The family requires every page to use the same spine edge and the two-edge
tail to start at a spine endpoint.  It does not cover a tail at a page
vertex (in particular Atlas 123), two-edge tails on an arbitrary chordal
core, or mixed maximal-clique sizes in general.

The graph, chromatic, target, Hilbert-square, moment, equality, and Atlas
identities are reconstructed exactly by
\begin{verbatim}
python codes/verify_triangle_book_two_edge_tail.py
\end{verbatim}


%%%%%%%%%%%%%%%%%%%%%%%%%%%%%%%%%%%%%%%%%%%%%%%%%%%%%%%%%%%%%%%%%%%%%%%%%%%%%%%
\section{What the Lean formalization actually proves}
%%%%%%%%%%%%%%%%%%%%%%%%%%%%%%%%%%%%%%%%%%%%%%%%%%%%%%%%%%%%%%%%%%%%%%%%%%%%%%%

The machine-checked development is
\texttt{lean/Taeyoung/Methods/BookTail/\{Core,Rows\}.lean}, ending at
\texttt{satisfiesLowerBound\_120}.  It follows this note step for step except at
one place, and the difference is worth stating precisely so that neither
document is read as a paraphrase of the other.

\paragraph{The departure: Lemma~\ref{lem:A-triangle} on $[1/2,1]$ only.}
This note proves
\[
 \int A\tau\geq(2p-1)\int Ad
 \qquad\text{for every }p\geq1/6,
\]
through the exact operator square \eqref{eq:operator-square}.  The formalization
proves the same inequality \emph{only for $p\geq1/2$}, and by a different
argument.  Write
\[
 M=\int A=\lVert d\rVert_2^2,\qquad
 N=\lVert A\rVert_2^2,\qquad
 B=\ip{A}{d}.
\]
Three applications of Cauchy--Schwarz give $B^2\leq NM$, $N\geq M^2$ and
$M\geq p^2$, and then the polynomial identity
\begin{equation}
 (2N-pM)^2-(2p-1)^2NM=(4N-M)(N-p^2M)
 \label{eq:lean-factorization}
\end{equation}
has both right-hand factors nonnegative: $N\geq M^2\geq p^2M$ gives the second,
and $M\geq p^2\geq1/4$ gives $4N\geq4M^2\geq M$ for the first.  With
$2N-pM\geq0$ and $2p-1\geq0$ this yields $2N-pM\geq(2p-1)B$, which combined
with \eqref{eq:first-reduction} is \eqref{eq:A-triangle}.

Both routes are correct; they differ in reach and in cost.  The operator square
is the stronger statement, valid on $[1/6,1]$, and it is the one to formalize if
that range is ever wanted --- it needs the six-term expansion of
$\int(A+ud+v)^2$ and the substitutions $\ip{h}{\one}=\lVert g\rVert_2^2$,
$A=pd+h$.  The factorization \eqref{eq:lean-factorization} needs no operator
language at all, and $p\geq1/2$ is exactly the interval on which the catalogue
proposition is asserted for a graph of chromatic number three.  No Atlas row
--- present or possible, since $\chi(R_m)=3$ for every $m$ --- requires the
extra range, so the formalization takes the cheaper route.

\paragraph{Scope: only $m\leq2$ is formalized.}
Theorem~\ref{thm:main} is proved here for every $m\geq1$, by Jensen for
$s\mapsto s^m$ under $\nu/B$.  The Lean file proves the case $m=2$, where that
Jensen step is a single Cauchy--Schwarz, and does not state the general $m$.
This loses nothing for the catalogue: $R_1$ is the five-vertex triangle with a
two-edge tail, already \texttt{verified} through the rooted-tree route, and
$R_3$ already has seven vertices.

\paragraph{Everything else transcribes directly.}
The pointwise estimate \eqref{eq:pointwise} is \texttt{Link.rootedTriangle\_ge},
which predates this note; $\int A=\int d^2$ is
\texttt{BookTail.integral\_pathOp}, one application of self-adjointness;
\eqref{eq:density} and \eqref{eq:first-moment} are
\texttt{BookTail.integral\_tailWeight} and
\texttt{integral\_tailWeight\_mul\_pageOp}; Lemma~\ref{lem:first-page} is
\texttt{BookTail.firstPage\_bound}, with the monotonicity of
$r\mapsto2r^2-pr$ replaced by the equivalent factorization
$2M^2-pM-p^3(2p-1)=(M-p^2)\bigl(2(M+p^2)-p\bigr)$.  The final division by $B$
is done as \texttt{le\_of\_mul\_le\_mul\_right}, so the argument never divides
by $2p-1$, which vanishes at the threshold.

\end{document}
